% Fermion TeX Engine — システム説明書（プロジェクト全体のしくみ）
% lualatex docs/system-manual.tex で組版
\documentclass[11pt,a4paper]{ltjsarticle}
\usepackage{luatexja}
\usepackage[a4paper,margin=22mm]{geometry}
\usepackage{amsmath}
\usepackage{booktabs}
\usepackage{array}
\usepackage{tikz}
\usetikzlibrary{positioning,arrows.meta}
\usepackage[hidelinks]{hyperref}

\newcommand{\code}[1]{\texttt{#1}}
\newcommand{\cs}[1]{\texttt{\textbackslash #1}}
\setlength{\parskip}{0.3\baselineskip}

\title{\textbf{Fermion TeX Engine システム説明書}\\[0.35em]
\large このプロジェクトは何を、どういう仕組みで「いじって」いるのか}
\author{TeX DOM Runtime プロジェクト}
\date{2026年7月3日}

\begin{document}
\maketitle

\begin{abstract}
Fermion TeX Engine は、組版システム LaTeX を「編集のたびに全体を作り直すバッチ処理」から
「文書を常駐させ、直した部分だけを数ミリ秒で組み直すランタイム」に作り替えたものです。
本書はプロジェクト全体の設計を一冊で説明します。特に力を入れたのは
\textbf{「何をいじって、何をいじっていないのか」}の完全な目録です。
本エンジンは LaTeX 本体（プログラム・組版アルゴリズム・フォント処理）には
1バイトも手を加えていません。外から注入しているのは、(1) OS の fork 機能を橋渡しする
71 行の C コード、(2) 組版結果を写し取る Lua プログラム、(3) 動作を変えずに記録だけを行う
約 40 行の TeX マクロ、の 3 種類だけです。この「注入物の少なさ」こそが、
プレビューが本物の PDF と一致することの根拠になっています。
\end{abstract}

\tableofcontents
\clearpage

%======================================================================
\section{この説明書について}

\paragraph{対象読者}
プログラミングの一般的な知識（プロセス、サーバー、ブラウザ程度の語彙）はあるが、
TeX の内部には詳しくない方を想定しています。TeX の予備知識は第
\ref{sec:tex-basics} 章で必要な分だけ説明します。

\paragraph{関連文書}
\begin{itemize}
\item \code{docs/01}〜\code{07}（Markdown）--- 実装の完全詳解。ソースコードと突き合わせて
  読むための技術資料で、開発中に踏んだ罠の記録（第6章）を含みます。
\item \code{docs/implementation-report.pdf} --- v3.2 の改良（プレビューと PDF の完全一致化）に
  絞ったやさしい読み物。
\item 本書 --- プロジェクト全体の仕組みを一冊で説明する「説明書」。上の二つの中間の詳しさです。
\end{itemize}

%======================================================================
\section{一言でいうと何をするソフトか}

LaTeX（ラテフ）は論文・技術書の定番組版システムです。仕上がりは印刷品質ですが、
\textbf{一文字直しただけでも文書全体を最初から処理し直す}ため、
結果を見るまで数秒〜数十秒かかります。

本エンジンはこれを次のように変えます。

\begin{center}
\begin{tabular}{p{0.30\textwidth}p{0.30\textwidth}p{0.30\textwidth}}
\toprule
& 従来の LaTeX & Fermion TeX Engine \\
\midrule
処理の形 & 使い捨てのバッチ処理 & \textbf{常駐プロセス（起動しっぱなし）} \\
編集 1 回のコスト & 文書全体の再処理 & \textbf{直した段落 1 つ分の組版（4〜14 ミリ秒）} \\
画面反映まで & 数秒〜 & \textbf{約 0.03 秒（打鍵ごと）} \\
見た目の計算 & LaTeX 本体 & \textbf{同じ LaTeX 本体}（無改変） \\
\bottomrule
\end{tabular}
\end{center}

つまり「Word のような即時性」と「LaTeX の印刷品質」の両立です。
脚注・図表の自動配置・目次・相互参照・日本語（luatexja）まで、
編集した瞬間にライブで更新されます。

%======================================================================
\section{予備知識 --- TeX はどうやって PDF を作るか}
\label{sec:tex-basics}

本エンジンが「どこに手を入れたか」を理解するには、
まず TeX が内部で何をしているかをざっくり知る必要があります。

\subsection{組版の流れ}

TeX が原稿（テキストファイル）から PDF を作るまでは、おおよそ次の段階を踏みます。

\begin{enumerate}
\item \textbf{読み取り}：原稿を 1 文字ずつ読み、マクロ（\cs{section} など）を展開する。
\item \textbf{段落の組版}：段落分の文字がたまると、有名な Knuth--Plass
  アルゴリズムが「どこで行を折り返すと最も美しいか」を段落全体で最適化する。
  ハイフネーション、カーニング（字間の微調整）、リガチャ（合字）もここで確定する。
\item \textbf{垂直リストへの追加}：出来上がった各行は「箱（ボックス）」として、
  \textbf{メイン垂直リスト}と呼ばれる縦一列の並びに追加されていく。
  箱と箱の間には「グルー（伸び縮みできる空白）」や「ペナルティ
  （ここで改ページすると減点いくつ、という印）」が挟まる。
\item \textbf{ページビルダー}：垂直リストがページの高さに達すると、
  ペナルティやグルーの情報から最良の改ページ位置を選び、1 ページ分を切り出す。
\item \textbf{出力ルーチン}：切り出された 1 ページ分に、脚注・図表（フロート）・
  ページ番号などを合成して、最終的なページの箱を作る（LaTeX ではここが
  \cs{output} ルーチンと呼ばれる仕組みで実装されている）。
\item \textbf{PDF 書き出し}：ページの箱の中の文字や線を PDF 命令に変換して書き出す。
\end{enumerate}

覚えておいてほしいキーワードは\textbf{「ノード」}です。
TeX の内部では、文字 1 つ、空白 1 つ、改ページ禁止の印 1 つが、
それぞれ「ノード」というデータとしてリスト状につながっています。
組版とは「ノードの列を作って寸法を確定させる」作業に他なりません。

\subsection{LuaTeX の特別なところ}

本プロジェクトが土台に選んだ \textbf{lualatex（LuaHBTeX）}は、TeX の現行エンジンのひとつで、
内部にスクリプト言語 \textbf{Lua} を積んでいます。重要なのは次の 2 点です。

\begin{itemize}
\item Lua のプログラムから、\textbf{組版済みのノードの列を直接読める}
  （どの文字がどの座標に置かれたか、グルーが実際に何 mm に伸びたか、まで見える）。
\item Lua から OS の共有ライブラリを読み込める（後述の C シムはこれで入れる）。
\end{itemize}

つまり LuaTeX は最初から「外から中身をのぞく公式の窓」を持っています。
本エンジンはこの窓だけを使い、\textbf{窓の内側（組版計算そのもの）には触れません}。

%======================================================================
\section{全体アーキテクチャ}

\subsection{登場するプロセス}

システムは 3 つの世界からできています：ブラウザ（画面）、Node.js のサーバー
（司令塔）、そして lualatex のプロセスの群れ（組版の実体）です。

\begin{center}
\begin{tikzpicture}[
  box/.style={draw, rounded corners, align=center, inner sep=6pt, font=\small},
  arr/.style={-{Stealth}, thick},
  node distance=9mm and 12mm
]
\node[box] (browser) {ブラウザ\\（エディタ＋プレビュー）};
\node[box, right=18mm of browser] (node) {Node.js サーバー\\（オーケストレータ）\\ \scriptsize server.js + engine-v3.js};
\node[box, below=of node, xshift=-32mm] (ckpt0) {ckpt0\\ \scriptsize lualatex 本体};
\node[box, right=8mm of ckpt0] (ckpt1) {ckpt1};
\node[box, right=8mm of ckpt1] (ckpt2) {ckpt2};
\node[right=6mm of ckpt2, font=\small] (dots) {$\cdots$};
\node[box, right=6mm of dots] (render) {レンダー子\\ \scriptsize 図を PDF 化して終了};
\draw[arr, <->] (browser) -- node[above, font=\scriptsize]{HTTP / SSE} (node);
\draw[arr, <->] (node) -- node[right, font=\scriptsize]{TCP（行指向プロトコル）} (ckpt0.north);
\draw[arr] (ckpt0) -- node[above, font=\scriptsize]{fork} (ckpt1);
\draw[arr] (ckpt1) -- node[above, font=\scriptsize]{fork} (ckpt2);
\draw[arr, <->] (node.south east) to[bend left=15] (render.north);
\end{tikzpicture}
\end{center}

\begin{description}
\item[ブラウザ] 左にテキストエディタ、右にプレビュー。編集すると差分だけをサーバーに
  送り、返ってきた「描画命令の差分」で該当ページだけを描き替えます。中身は徹底して
  薄く、組版の知識は持ちません。
\item[Node.js サーバー（オーケストレータ）] 文書の「真実」を握る側。ソース全文、
  ブロック分割、ラベル表、ページ、フォント台帳を保持し、lualatex プロセス群に
  仕事を割り振ります。ページへの割り付け（第 \ref{sec:pagebuilder} 章）もここで行います。
\item[lualatex プロセスの群れ] 起動時に 1 個だけ本物の lualatex を立ち上げ、
  あとは fork（自己複製）で増えます。各プロセスが「文書を途中まで処理した状態」の
  生きたセーブポイントです（第 \ref{sec:checkpoint} 章）。
\end{description}

\subsection{責任分担の原則}

\begin{itemize}
\item \textbf{文字の並び・行の折り返し・あらゆる寸法の決定} --- lualatex（無改変の本物）。
\item \textbf{どのブロックを組み直すかの判断、ページへの割り付け、差分配信} --- Node.js。
  ただし割り付けは TeX の規則書をそのまま移植したもので、独自の判断はしません。
\item \textbf{描画} --- ブラウザ。ただし TeX が使ったフォントファイルそのものを使い、
  ブラウザ独自の文字詰め機能はすべて無効化します（第 \ref{sec:render} 章）。
\end{itemize}

%======================================================================
\section{何を「いじって」いるのか --- 注入物の完全目録}
\label{sec:injections}

本プロジェクトの正確な要約は次の一文です。

\begin{quote}
\textbf{「lualatex が計算する内容」は一切変えず、「いつ・どの単位で計算するか」と
「結果をどこへ出すか」だけを差し替えた。}
\end{quote}

\subsection{いじって「いない」もの}

先に、無改変のものを列挙します。ここが本エンジンの信頼性の土台です。

\begin{center}
\begin{tabular}{p{0.42\textwidth}p{0.52\textwidth}}
\toprule
無改変のもの & 意味 \\
\midrule
lualatex（LuaHBTeX）のバイナリと C ソース & TeX Live のものをそのまま起動するだけ \\
Knuth--Plass 行分割・数式組版・ハイフネーション & 段落や数式の見た目を決める計算はすべて原作者のコードが実行 \\
フォント処理（luaotfload / HarfBuzz） & リガチャ・カーニング・字形選択は従来どおりエンジン内で確定 \\
LaTeX カーネルとフォーマット & \code{latex.ltx} 無改変、\code{lualatex.fmt} をそのままロード \\
パッケージ群（amsmath, tikz, booktabs, luatexja, \ldots） & 無改変 \\
PDF エクスポート & 「PDF を保存」は素の lualatex の 2 回実行。ライブ表示との一致検算にもなる \\
\bottomrule
\end{tabular}
\end{center}

\subsection{注入物 1：C シム \code{tdomfork.c}（71 行）}

LuaTeX の Lua には、プロセスを複製する \code{fork} を呼ぶ標準手段がありません
（同梱の FFI ライブラリは Apple Silicon 非対応）。そこで
\code{fork} / \code{waitpid} / \code{\_exit} / \code{getpid} /
「子プロセスの後始末を OS に任せる設定」の 5 つの OS 機能を Lua に橋渡しするだけの
共有ライブラリを C で書きました。全 71 行で、組版には一切関与しません。
初回起動時にオーケストレータが \code{cc} で自動コンパイルします。

\subsection{注入物 2：Lua デーモン \code{daemon.lua}（約 900 行）}

lualatex の中に読み込まれる Lua プログラムで、仕事は 4 つです。

\begin{enumerate}
\item \textbf{通信}：オーケストレータと TCP ソケットでつながり、命令
  （「このブロックを組め」など）を待ち受ける。
\item \textbf{凍結}：命令待ちの間、プロセスを眠らせる（＝セーブポイントとして保存）。
\item \textbf{収穫}：組版済みのノード列を歩いて、「どのフォントのどの文字がどの座標にあるか」
  「グルーやペナルティの実寸」を JSON にしてオーケストレータへ送る。
\item \textbf{出荷}：図などを本物の PDF ページとして書き出す（精密レンダー用）。
\end{enumerate}

すべて LuaTeX の公式 API の範囲内で、ノードを「読む」ことはしても
組版結果を「書き換える」ことはしません。

\subsection{注入物 3：TeX マクロのシム（約 40 行）}

LaTeX の一部の命令を、「\textbf{元の動作をそのまま実行した上で}、ついでに Lua へ
記録する」形に包み直します。増分更新に必要な依存関係（どのブロックがどのラベルを
参照しているか等）を知るためです。主なものを挙げます。

\begin{center}
\small
\begin{tabular}{lll}
\toprule
対象 & 目的 & 組版への影響 \\
\midrule
\cs{label}, \cs{ltx@label} & ラベル値の記録と即時定義 & なし \\
\cs{ref}, \cs{pageref}, \cs{eqref} & 「誰がどのラベルを参照するか」の記録 & なし \\
\cs{cite}, \cs{@bibitem} 系 & 引用と文献番号の記録・即時解決 & なし \\
\code{figure}/\code{table} 環境 & フロートの中身と配置指定の捕捉 & ライブ層のみ（印刷は無影響） \\
\cs{@starttoc} & 目次を当方管理のファイルから読ませる & なし \\
ページ設定（余白原点ほか） & 図の切り出しの原点合わせ & ライブ層のみ \\
フォント準備・数式代替フォント計測 & 起動時に一度だけ実行 & なし（計測用の箱は捨てる） \\
\bottomrule
\end{tabular}
\end{center}

フロート環境の包み直しも、箱の作り方（幅・字下げ復元・フォントリセット）は
LaTeX 本体の \cs{@xfloat} と同一手順にしてあり、捕捉される箱は
本来の出力ルーチンが受け取るはずだった箱と同じものです。

\subsection{注入物 4：ページビルダーの「休眠」設定}

TeX 内蔵のページビルダー（改ページ係）は、常駐運転では眠らせておきます。
改ページはオーケストレータ側で行うためです（理由は第 \ref{sec:pagebuilder} 章）。
眠らせ方は 3 つの設定の組み合わせです：
「ページの高さを事実上無限大にする（ページが決して満杯にならない）」、
「脚注の材料を横取りさせず行の位置に留める」、
「ページの先頭に番兵の空箱を 1 つ置く（TeX が『ページはまだ空だ』と誤解して
先頭の空白を捨てたり余計な空白を足したりするのを防ぐ）」。
ユーザー原稿に \cs{newpage} などの強制改ページがあった場合に備えて、
発火した出力ルーチンの中身をそっくり戻して「ここで改ページ」という印だけを残す
安全網も入れてあります。

\subsection{安全性について --- \code{-{}-shell-escape} の意味}

C シムを Lua に読み込ませるため、lualatex は \code{-{}-shell-escape} 付きで起動します。
これは「TeX 文書が外部プログラムを起動してよい」という許可であり、
\textbf{信頼できる自分の文書を自分のマシンで組む}ためのローカル開発ツールとしての
前提です。見知らぬ他人の \code{.tex} をそのまま流し込む用途は想定していません。

%======================================================================
\section{心臓部 1 --- fork チェックポイント}
\label{sec:checkpoint}

\subsection{fork は「TeX の完全なセーブデータ」である}

TeX の内部状態は巨大です：数万のマクロ定義、文字種の表、ロード済みフォント、
章や数式の番号カウンタ……。これをアプリの力で保存・復元するのは絶望的ですが、
OS から見れば\textbf{ただのプロセスメモリ}です。

POSIX の \code{fork(2)} は動いているプロセスの完全な複製を作るシステムコールで、
現代の OS では\textbf{コピーオンライト}（実際に書き換わったメモリページだけを
あとから複製する方式）で実装されています。だから：

\begin{itemize}
\item fork 自体は一瞬（実測 0.2 ミリ秒）。
\item 複製のメモリ代は「その後書き換わった分」だけ。
\item \textbf{fork した親を眠らせて残せば、それが「その瞬間の TeX 完全状態の
  生きたセーブポイント」になる。}
\end{itemize}

文書をブロック（≒段落・見出し・数式環境。空行で区切る）ごとに組版しながら、
各境界で fork して親を残すと、次のようなプロセスの鎖ができます。

\begin{center}
\code{ckpt0}（前書き処理済み）→ \code{ckpt1}（ブロック1まで済み）→
\code{ckpt2} → $\cdots$ → \code{ckptN}（文書末尾まで済み）
\end{center}

ユーザーがブロック $K$ を編集したら、(1) \code{ckpt}$_K$ 以降（古い続きに基づく
状態）を破棄し、(2) \code{ckpt}$_{K-1}$ に「新しいブロック $K$ を組め」と命じます。
子が組版して新しい \code{ckpt}$_K$ になります。プロセス起動もフォーマット読み込みも
フォント読み込みも発生しません。これが数ミリ秒応答の正体です。

\subsection{凍結のからくり}

「プロセスを眠らせてセーブポイントにする」は、たった 1 行の TeX マクロと
Lua の受信待ちで実現されています。TeX は
\cs{directlua}\code{\{tdom\_wait()\}} の内部で永久に停止し（ソケットの受信待ちなので
CPU は消費しません）、fork で生まれた子だけが Lua から戻って組版を続けます。
組版と報告を終えた子は、再び同じ待ち受けに入って次のセーブポイントになります。
1 本のマクロループの上を、fork の世代が次々に流れていく構造です。

\subsection{編集 1 回で起きること}

\begin{enumerate}
\item ブラウザから「何文字目から何文字目をこう置き換えた」という差分が届く。
\item オーケストレータが本文をブロックに分け直し、内容ハッシュで前回と突き合わせて
  「変わったブロック」を特定する（1 ミリ秒未満）。
\item 変わったブロック以降のセーブポイントを破棄し、直前のセーブポイントに
  組版を命じる（fork → 組版 → 結果の JSON → 新セーブポイント昇格）。
\item \textbf{収束の自己検証}：続きのブロックも 1 つ組み直してみて、出力と状態
  （カウンタ類の出口値）が前回と完全一致したら、そこで止めてよいと判断する。
  一致しなければ（式番号がずれた等）一致するまで先へ進む。
  「LaTeX をもう一回コンパイル」をブロック単位に局所化した形で、
  \textbf{停止判定そのものが検算}になっています。
\item ページ割り付けを再計算し、変わったページの描画命令だけをブラウザへ送る。
\end{enumerate}

なお v3.2 からは、状態ベクトルに「直前の行の深さ」「見出し直後かどうか」という
\textbf{レイアウト上の連続状態}も含まれており、段落間の空白に影響する変化も
この自己検証の網に掛かります。

\subsection{プロセス数の管理}

セーブポイント＝プロセスなので、巨大文書では数を絞ります（既定で最大 64 個）。
間引かれた区間の編集は、少し手前のセーブポイントから数ブロック余分に組み直すだけです
（1 ブロック数ミリ秒なので体感差はありません）。

%======================================================================
\section{心臓部 2 --- 組版結果の収穫}
\label{sec:harvest}

\subsection{本流の上で組み、写し取る}

各ブロックは、TeX が文書を組み上げていく\textbf{本物のメイン垂直リストの上で}
組版されます（第 \ref{sec:injections} 章の「休眠」設定により、ページビルダーが
横から切り取らないようにした上で）。組み終わったら、Lua デーモンがたまった
ノード列を丸ごと写し取り、リストを空に戻してから凍結します。

この方式の要点は、\textbf{段落と段落の間の情報も全部本物になる}ことです。
段落間の空白は「直前の行の深さ」に依存し、見出し直後の段落は字下げをしない——
こうした前後関係は、プロセスのセーブポイント連鎖の中で自然に引き継がれます。
実験により、写し取ったノード列は「素の lualatex が文書全体を一気に組んだ場合」と
\textbf{バイト単位で一致}することを確認済みです。

\subsection{run 分割の契約 --- 座標を全部送らずに済む理由}

写し取る情報は「行の箱」の中のグリフ（文字の字形）たちです。
1 文字ずつ座標を送るとデータが膨れるため、次の契約を使います。

\begin{quote}
グリフの連なり（run）は、カーニングや空白（kern / glue のノード）が現れるたびに
必ず分割する。すると run の内部にはもう位置調整が存在しないので、
\textbf{「先頭の $x$ 座標」だけ送れば、あとはフォントの字送り幅を足すだけで
TeX と同じ座標が再現できる}。
\end{quote}

ブラウザ側は TeX が使ったのと同じフォントファイルで描くので（次章）、
字送り幅の出どころも同一です。これでペイロードが約 1 桁縮みます。

\subsection{収穫される JSON の中身}

ブロック 1 つの収穫結果には、行の箱（高さ・深さ・run の列）、実寸グルー
（自然長と伸縮成分）、ペナルティ、脚注の中身（行位置に紐づく）、フロートの実体
（箱の寸法と中身、配置指定）、使用フォントの実ファイルパス、記録されたラベルと参照、
カウンタの出口値、が入っています。オーケストレータはこれを材料に
ページ割り付けと差分判定を行います。

%======================================================================
\section{心臓部 3 --- ページ割り付け（規則書の移植）}
\label{sec:pagebuilder}

\subsection{なぜ Node 側で改ページするのか}

増分更新の要は「編集していないページを触らない」ことです。改ページを TeX に任せると、
1 ページ分の材料が確定するまで TeX を走らせる必要があり、ブロック単位のセーブポイントと
相性が悪くなります。そこで\textbf{改ページとページ合成だけ}をオーケストレータ側で
行います。

ただし我流の割り付けでは本物とずれます。そこで v3.2 で、この部分を
\textbf{TeX の原典プログラム（改ページの数式まで含む）と LaTeX の出力ルーチンの
逐語的な移植}に置き換えました。具体的には：

\begin{itemize}
\item \textbf{改ページ判定}：TeX と同じ「改ページしてよい場所」の定義、同じコスト計算
  （空きの悪さ badness ＋ ペナルティ。badness は原典と同じ整数演算）、
  同じ「最良候補を覚えておき、あふれた時点で確定」の手順。
\item \textbf{ページ上端の空白}（\cs{topskip}）、\textbf{行の深さの繰り入れ}
  （\cs{maxdepth}）、\textbf{脚注の場所の先取り会計}（脚注が増えるほど本文の
  取り分が減る計算）も原典どおり。
\item \textbf{フロート配置}：図表を「今のページの上か・下か・その場か・後回しか・
  専用ページか」に振り分ける LaTeX の判定式（回数制限・面積割合・順序保存の
  すべて）を条件式ごと移植。
\item \textbf{ページ合成}：本文＋脚注の区切り線（クラス定義を実測した寸法で再生）＋
  脚注＋上下フロートを、TeX と同じグルー分配（伸縮の次数まで考慮）で最終座標にする。
\end{itemize}

このファイルの中に\textbf{発明された寸法は 1 つもありません}。
すべての定数は、起動時に本物の LaTeX から伸縮成分込みで実測して受け取ります。

%======================================================================
\section{描画層 --- ブラウザで「印刷と同じ絵」を出す}
\label{sec:render}

\subsection{フォントは TeX が使った実ファイルを配る}

プレビューの文字は、TeX Live のディスクにあるフォントファイル
（Latin Modern、原ノ味フォント等）\textbf{そのもの}をサーバーが配信し、
ブラウザの \code{@font-face} として登録して描きます。
TeX と同じ字形・同じ字送り幅が使われるので、run 契約（第 \ref{sec:harvest} 章）が
成立します。

\subsection{ブラウザの整形機能をすべて切る}

ブラウザは放っておくと独自にカーニングやリガチャを再適用し、TeX の結果とずれます。
そこで描画にはカーニング無効・リガチャ無効・字間調整ゼロを指定します。
リガチャは TeX がすでに実行済みで、文字列には合字そのものが入っています。

\subsection{数式の Type1 問題と「双子」フォント}

古典的な LaTeX 数式のフォント（Computer Modern）は 1980 年代の Type1 形式で、
ブラウザは表示できません。そこで公式後継の OpenType 版（Latin Modern）を
「双子」として使い、字形番号の対応表で写像します。それでも寸法が合わない
大型記号（大きな積分記号など）を含むブロックは、次の精密レンダー行きにします。

\subsection{2 層戦略 --- 即時グリフ＋精密チャンク}

TikZ の図やグラフのように「グリフでは描けないもの」（生の PDF 描画命令）は、
2 段構えで表示します。

\begin{enumerate}
\item \textbf{即時層}：打鍵直後は、まず文字ベースの近似を表示（体感ゼロ遅延のため）。
\item \textbf{精密層}：裏でレンダー子プロセスが該当ブロックを\textbf{本物の PDF ページ}
  として書き出し、SVG 画像に変換してプレビューに貼り込む（1〜2 秒後に自動置換）。
  貼られる絵は本物の PDF の絵そのものなので、\textbf{原理的に印刷と同じ}見た目です。
\end{enumerate}

レンダー子はセーブポイントから fork されるため、文書の途中の状態
（カウンタ・フォント・定義）を完全に引き継いだうえで図だけを出荷できます。

%======================================================================
\section{打鍵から画面まで --- 時系列}

キーを 1 つ打ってから画面が変わるまでの実測の内訳です（英語デモ文書の場合）。

\begin{center}
\begin{tabular}{lll}
\toprule
段階 & 担当 & 時間の目安 \\
\midrule
差分の送信 & ブラウザ & 打鍵ごと即時 \\
ブロック分割と差分特定 & Node & 1 ミリ秒未満 \\
fork ＋ 変更ブロックの組版 & lualatex & 4〜14 ミリ秒 \\
収穫（ノード走査 → JSON） & Lua デーモン & 数ミリ秒 \\
ページ割り付け・描画命令の差分化 & Node & 1 ミリ秒前後 \\
該当ページの描き替え & ブラウザ & 数ミリ秒 \\
\midrule
合計（打鍵 → 画面） & & \textbf{約 30 ミリ秒} \\
\bottomrule
\end{tabular}
\end{center}

式を 1 本挿入して以降の式番号が全部ずれるような編集でも、影響したブロックだけの
再組版で約 50 ミリ秒です。日本語文書（luatexja）は 1 編集あたり実測 5 ミリ秒前後。
TikZ の図の編集は即時近似が数十ミリ秒、精密画像への置換が 1〜2 秒です。

%======================================================================
\section{正しさの保証}

「プレビュー＝本物の PDF」を主張だけで終わらせないため、二重の仕組みがあります。

\paragraph{（1）収束の自己検証（常時・毎編集）}
第 \ref{sec:checkpoint} 章のとおり、増分更新の停止判定そのものが
「もう 1 ブロック組んで前回と完全一致するか」という検算です。

\paragraph{（2）答え合わせ機（開発時・リリース時）}
\code{tools/verify-layout.mjs} は、同じ原稿を本エンジンと素の lualatex の両方で
処理し、\textbf{全ページ・全行の位置}を数値で突き合わせます。真値は PDF ファイルの
中身（描画命令）を直接読んで取り出します。現在、同梱の全テンプレート・デモ文書・
日本語実文書で\textbf{全行が 0.02 bp（約 7 ミクロン ＝ 髪の毛の 1/10 以下）以内で一致}し、
編集を繰り返した後の状態でも一致します（\code{tools/verify-edits.mjs}）。
統合テストは 30 本すべて通過しています。

%======================================================================
\section{コードマップと API}

\subsection{リポジトリの地図}

\begin{center}
\small
\begin{tabular}{llp{0.47\textwidth}}
\toprule
ファイル & 行数 & 役割 \\
\midrule
\code{server.js} & 269 & HTTP/SSE サーバー。バックエンド選択、編集受付、配信 \\
\code{engine/checkpoint/engine-v3.js} & 約 1{,}700 & オーケストレータ本体。差分・依存・収束・表示リスト \\
\code{engine/checkpoint/daemon.lua} & 約 900 & TeX 内デーモン。凍結・収穫・出荷 \\
\code{engine/checkpoint/pagebuilder.js} & 約 830 & ページ割り付け（TeX 規則書の移植） \\
\code{engine/checkpoint/tdomfork.c} & 71 & fork の橋渡し C シム \\
\code{engine/checkpoint/mathmap.js} & 186 & 数式レガシーフォント → 双子 OTF の対応表 \\
\code{engine/segmenter.js} ほか & --- & ブロック分割・ハッシュ差分（全世代共有） \\
\code{engine/engine-lua.js}, \code{engine/luatex/} & --- & v1 バックエンド（比較用の参照実装） \\
\code{engine/engine.js} ほか & --- & v0 内部エンジン（TeX 不要環境のフォールバック） \\
\code{web/} & 約 1{,}000 & ブラウザ側（エディタ・プレビュー・インスペクタ） \\
\code{tools/} & 約 580 & 答え合わせ機・デバッグ道具 \\
\code{tests/} & --- & 統合テスト 30 本 \\
\code{docs/} & --- & 技術詳解 7 章＋本書＋レポート \\
\bottomrule
\end{tabular}
\end{center}

\subsection{3 世代のバックエンド}

リポジトリには設計の進化が 3 つのバックエンドとして残っており、起動時に選べます
（環境変数 \code{TDOM\_BACKEND}）。

\begin{itemize}
\item \textbf{v0 internal}：TeX 不要の自前 JS 近似。アーキテクチャの参照実装兼
  フォールバック。
\item \textbf{v1 lualatex}：編集のたびに変更ブロックだけを新しい lualatex プロセスで
  組む方式（約 0.4 秒/編集）。fork が使えない環境の代替。
\item \textbf{v3 checkpoint（既定）}：本書で説明した常駐 fork 方式。
\end{itemize}

\subsection{HTTP API}

\begin{center}
\small
\begin{tabular}{lp{0.6\textwidth}}
\toprule
ルート & 内容 \\
\midrule
\code{GET /doc} & 初期ロード一式（ソース・全ページ描画命令・ジオメトリ・フォント鍵） \\
\code{POST /edit} & 範囲編集 \code{\{start,end,text\}} → 差分レポート＋ページパッチ \\
\code{POST /open} & 文書の開き直し（本文指定 or テンプレート指定） \\
\code{GET /events} & SSE。他ウィンドウ同期・図の精密画像の非同期差し替え \\
\code{GET /font/…} & TeX が実際に使ったフォントファイル \\
\code{GET /chunk/…} & 図などの精密 SVG \\
\code{GET /pdf} & 素の lualatex 2 回実行による正式 PDF \\
\code{GET /dom} & ブロック・依存・ページ対応の観測用 JSON \\
\code{GET /templates} & テンプレート一覧 \\
\bottomrule
\end{tabular}
\end{center}

%======================================================================
\section{動かし方と動作条件}

\begin{itemize}
\item 必要なもの：Node.js、TeX Live（\code{lualatex}）、poppler（\code{pdftocairo}）、
  C コンパイラ（\code{cc}）。npm の外部依存は\textbf{ゼロ}です。
\item 起動：リポジトリで \code{npm start} → ブラウザで \code{http://127.0.0.1:4633}。
\item OS：fork(2) を使うため macOS / Linux（Windows は WSL で可）。
\item バックエンド切替：\code{TDOM\_BACKEND=internal|lualatex|checkpoint}。
\end{itemize}

\paragraph{既知の未対応}
脚注のページまたぎ分割、二段組、欄外注（marginpar）、
\code{figure}/\code{table} 以外の独自フロート環境。いずれも検出時に診断を出し、
だんまりで近似することはしません。

%======================================================================
\section{用語集}

\begin{description}
\item[組版（くみはん）] 文字・図をどの位置にどんな間隔で置くかを決める作業。
\item[ノード] TeX 内部のデータ単位。文字・空白・罫線・ペナルティなどが
  リスト状につながる。
\item[メイン垂直リスト] 組み上がった行の箱が縦に積まれていく、文書の「本流」。
\item[グルー] 伸び縮みできる空白。自然長・伸び量・縮み量を持つ。
\item[ペナルティ] 「ここで改ページ（改行）すると減点いくつ」という印。
  負なら「むしろここで切れ」。
\item[ページビルダー / 出力ルーチン] 垂直リストからページを切り出し、
  脚注や図を合成する TeX / LaTeX の仕組み。本エンジンでは休眠させ、
  移植版がオーケストレータ側で同じ計算を行う。
\item[フロート] 図表を「近くの良い場所」に自動配置する LaTeX の仕組み。
\item[fork / コピーオンライト] プロセスの完全な複製を一瞬で作る OS 機能と、
  その実装方式。本エンジンのセーブポイントの実体。
\item[チェックポイント] 「ブロック $k$ まで処理し終えた TeX の完全状態」を
  保持したまま眠っているプロセス。
\item[ガレー] 収穫された組版結果（行の箱と run の列）。活版印刷の「ゲラ」に由来。
\item[run] 途中に位置調整を含まないグリフの連なり。先頭座標だけで再現できる。
\item[チャンク] 図などを本物の PDF から切り出した精密 SVG 画像。
\item[bp] 長さの単位（ビッグポイント）。1 bp ≒ 0.353 mm。PDF の内部単位。
\end{description}

\vspace{1.5em}
\begin{center}
\rule{0.6\textwidth}{0.4pt}\\[0.5em]
\small 本書自体も、本エンジンと同じ lualatex（TeX Live）で組版されています。
\end{center}

\end{document}
